\documentclass[11pt]{article}

\usepackage{series}
\usepackage{amsmath,amssymb,amsfonts}
\usepackage{geometry}
\usepackage{hyperref}
\usepackage{enumitem}

\geometry{margin=1in}

% ---------- metadata ----------
\title{The Modal Triplet Theory Program A2:\\
Conditional Computability and Finite Prediction Depth}
\author{Peter Nero}
\date{July 2026}

\begin{document}

\maketitle

\begin{abstract}
We separate four notions that were conflated in the first version of Program
A2: finite admissible prediction depth, underdetermination at a branching
relation, algorithmic undecidability, and computational lower bounds.
Prediction depth is a chart- and protocol-relative iteration diagnostic.  It
is finite only under a certified finite barrier along the declared
continuations; finite depth alone does not imply undecidability.  A selection
reachability problem becomes algorithmically undecidable if a selected MTT
realization robustly and uniformly embeds a universal two-counter machine,
preserves every finite machine prefix, and identifies halting with a selected
target event.  We prove this conditional transfer theorem and also prove the
complementary finite-capacity result: a fixed finite carrier with computable
transitions has decidable reachability and therefore cannot satisfy the
unbounded embedding contract.  Neither non-Markovianity, noninjective
projection, branching, nor nil termination supplies that contract by itself.
Computational irreducibility additionally requires a specified computation
model and a lower-bound theorem.  Probability, quantization, physical time,
and universal limits of control remain separate source obligations.
\end{abstract}

\section*{Revision note for version 2}
\addcontentsline{toc}{section}{Revision note for version 2}
\paragraph{Supersedes.}
Version~1 of Program A2.
\paragraph{Reason.}
The former paper used ``undecidable'' for several different failures of
prediction and inferred algorithmic undecidability from finite admissibility,
branching, and nil termination without a declared input language or reduction.
It also combined unbounded counter storage with a finite admissibility horizon.
\paragraph{Resolution.}
Version~2 introduces a computable reachability presentation, states the exact
robust two-counter embedding contract, proves undecidability only conditional
on that contract, and proves that a fixed finite carrier cannot realize it.
It also gives correctly scoped prediction-depth and control statements.
\paragraph{Retained content.}
Prediction as admissible continuation, protocol-relative prediction depth,
and the distinction between local prediction and a global effective
description remain useful.
\paragraph{Open boundary.}
No selected MTT realization is presently proved to implement unbounded,
robust counter storage and operations.  No complexity lower bound, Born
measure, quantization theorem, physical time variable, or universal controller
no-go follows from A2 alone.

\tableofcontents
\newpage

\section{Scope and Logical Separation}

Program A0 supplies typed upper evolution, reductions, admissible domains, and
an iteration-count prediction-depth diagnostic~\cite{A0}.  Program A1 supplies
chart-persistence kinematics and only a conditional bridge to physical
spacetime propagation~\cite{A1}.  The present paper adds a computability layer.
It does not turn descriptive order into physical time or a chart obstruction
into a theorem of mathematical logic.

\subsection{Four different questions}

Let a finite string describe an initial state and a protocol.  The following
questions are distinct.
\begin{enumerate}[label=(\roman*)]
\item \emph{Semantic determinacy:} does the selected data define one
continuation, several continuations, or no continuation?
\item \emph{Decidability:} is there an algorithm that returns the correct
yes/no answer on every input in a declared language?
\item \emph{Complexity:} among algorithms that decide the language, what
resources are required as a function of input length?
\item \emph{Simulation cost:} for a specified trajectory, how much work is
required to reproduce a chosen number of steps?
\end{enumerate}
Branching may make an actual branch underdetermined while existential or
universal reachability remains decidable.  A long simulation may be expensive
without being undecidable.  Conversely, undecidability is a uniform statement
over an infinite input family, not a claim that every instance is hard.

\begin{remark}[Terminology]
In this paper, ``undecidable'' has its standard algorithmic meaning.  A question
outside the domain of an encoding is \emph{undefined} there.  A question with
several admissible answers is \emph{underdetermined} until a selector is added.
Neither condition is renamed undecidability.
\end{remark}

\section{Computable Admissible Presentations}

Computability requires more structure than a topological atlas.

\begin{definition}[Computable admissible presentation]\label{def:presentation}
A computable admissible presentation is a tuple
\[
 \mathfrak C=(\Sigma,S,\iota,F,A,H)
\]
with the following data.
\begin{itemize}
\item $\Sigma$ is a finite alphabet and $S\subseteq\Sigma^\ast$ is a decidable
set of finite state descriptions, equipped with a computable
rational-valued metric $d_S$ when robustness is claimed.
\item $\iota:\Sigma^\ast\to S$ is a total computable input encoder on the
declared input language.
\item $F:S\rightharpoonup S$ is a partial computable one-step map for one fixed
protocol, with decidable domain.
\item $A\subseteq S$ is a decidable admissible set.
\item $H\subseteq A$ is a decidable target or selection-event set.
\end{itemize}
An orbit is followed only while $F$ is defined and remains in $A$.
\end{definition}

The deterministic definition is sufficient for the reduction below.  A
branching system instead uses a computably presented relation
$R\subseteq S\times S$ and must declare whether reachability means
``there exists a branch'' or ``every branch.''

\begin{definition}[Selection reachability language]\label{def:sel}
For a presentation $\mathfrak C$, define
\[
\operatorname{SEL}(\mathfrak C)
:=
\left\{
w\in\Sigma^\ast:
\begin{array}{l}
\text{for some }n\geq0,\ F^k(\iota(w))\text{ is defined and in }A\\
\text{for }0\leq k\leq n,\text{ and }F^n(\iota(w))\in H
\end{array}
\right\}.
\]
\end{definition}

Because a target hit has a finite witness, this language is recursively
enumerable under Definition~\ref{def:presentation}.  Decidability requires
more: a decision procedure must also halt on every nonmember.

\subsection{Finite recognizable horizons}

\begin{proposition}[Decidability with recognizable termination]
\label{prop:recognizable}
Suppose every orbit of $\mathfrak C$ reaches $H$, leaves $A$, or repeats a
previous state after finitely many steps, and each of these events is
recognizable from the finite state description.  Then
$\operatorname{SEL}(\mathfrak C)$ is decidable.
\end{proposition}

\begin{proof}
Simulate the orbit while storing visited states.  Accept on entering $H$.
Reject on leaving $A$, on an undefined transition, or on the first repeated
state.  By hypothesis one of these cases occurs after finitely many steps.
\end{proof}

\begin{remark}
A finite horizon is not enough if its endpoint is only a semantic statement
that the chosen encoding no longer represents the intended upper object and
there is no effective test for that loss.  Conversely, an instance-dependent
horizon need not be known in advance if exit is effectively recognizable.
The algorithm may simply simulate until the recognizable exit.
\end{remark}

\section{Admissible Prediction Depth}

We now retain the useful A2 diagnostic in a form compatible with A0.

\begin{definition}[Protocol-relative prediction depth]
Let $T_{A,\pi}:P_0(A)\rightharpoonup P_0(A)$ be a declared reduced step map for
an admissible domain $A$ and protocol $\pi$.  For $x\in A$, set
\[
N_{\max}(x;A,\pi)
:=
\sup\left\{
n\in\mathbb N:
T_{A,\pi}^{k}(P_0x)\text{ is defined and in }P_0(A)
\text{ for }0\leq k\leq n
\right\}.
\]
\end{definition}

This number counts iterations.  A conversion to elapsed time, proper time,
distance, energy, or an information budget requires a realization-specific
bridge.

\begin{theorem}[Certified finite depth]\label{thm:depth}
If there is an integer $d<\infty$ such that every declared continuation from
$P_0x$ either leaves $P_0(A)$, becomes undefined, or reaches a separately
declared branch ambiguity by step $d$, then
\[
N_{\max}(x;A,\pi)\leq d
\]
for the corresponding continuation convention.
\end{theorem}

\begin{proof}
No continuation satisfying the defining conditions can contain more than
$d$ admissible, unambiguous steps.
\end{proof}

\begin{remark}[What does not prove finite depth]
The mere existence of a selection front or nil boundary somewhere in a
reachable graph does not give a uniform bound on every path.  One must prove
that the declared paths meet the barrier in finite depth, or supply a
coercive margin or ranking function that yields such a bound.
\end{remark}

\begin{proposition}[Refinement monotonicity under a projection contract]
\label{prop:refinement}
Let a refined transition system project stepwise to a coarse one, let every
refined admissible path project to a coarse admissible path, and let the
refined determinacy criterion imply the coarse criterion.  Then refined
prediction depth cannot exceed coarse prediction depth.
\end{proposition}

\begin{proof}
Every refined path counted at depth $n$ produces a coarse path counted at
depth $n$.  Taking suprema gives the claim.
\end{proof}

Without this path-projection contract, refinement can expose an obstruction,
remove a spurious obstruction, or change the protocol, so monotonicity is not
automatic.

\section{Branching, Nil Boundaries, and Algorithms}

\begin{proposition}[Branching is not undecidability]\label{prop:branch}
Suppose a computable, finitely branching relation $R$ returns an effective
finite successor list for each state.  For a fixed finite depth $d$, both
\[
\exists\text{ a branch reaching }H\text{ by }d
\quad\text{and}\quad
\forall\text{ branches, }H\text{ is reached by }d
\]
are decidable.
\end{proposition}

\begin{proof}
Enumerate the finite computation tree to depth $d$ and inspect its leaves.
\end{proof}

Thus a branch selector may be absent even though reachability questions are
algorithmically decidable.  Conversely, an infinite computable transition
system can have an undecidable reachability language even when each state has
only one successor.

\begin{proposition}[Nil boundary semantics]\label{prop:nil}
If a chart continuation is undefined beyond a certified nil boundary, then a
question whose semantics requires that continuation is undefined in that
chart system.  It becomes a yes/no decision problem only after a total
extension, rejection convention, or alternative chart relation is specified.
\end{proposition}

\begin{proof}
A partial map has no value outside its domain.  Algorithmic decidability is a
property of a total Boolean-valued decision problem, so the missing semantic
value must first be completed by declared data.
\end{proof}

\section{Two-Counter Machines and the Missing MTT Contract}

A deterministic two-counter machine has a finite control set $Q$, two
registers in $\mathbb N$, and instructions
\textsc{inc}, \textsc{decjz}, and \textsc{halt}.  A fixed universal machine
can receive its program and input through the initial register values.
Its halting language is undecidable~\cite{Minsky,Turing}.

\begin{definition}[Robust uniform MTT embedding]\label{def:embedding}
Let $\mathcal M$ be a fixed universal two-counter machine with configuration
set
\[
\operatorname{Conf}(\mathcal M)=Q\times\mathbb N^2
\]
and partial step map $\tau$.  A robust uniform embedding of $\mathcal M$ into
an MTT presentation consists of:
\begin{enumerate}[label=(E\arabic*)]
\item a computable injective code
$J:\operatorname{Conf}(\mathcal M)\to A$ and a total computable reduction map
$g$ from machine inputs to presentation inputs such that
\[
\iota(g(w))=J(c_0(w));
\]
\item \emph{step fidelity},
\[
F(J(c))=J(\tau(c))
\]
for every nonhalting configuration on every encoded run;
\item \emph{unbounded admissible persistence}: every finite prefix of every
encoded machine run remains in the domain of $F$ and in $A$;
\item \emph{target fidelity}: $J(c)\in H$ if and only if $c$ is a halting
configuration;
\item \emph{uniformity}: $F,A,H$, and the decoding convention are fixed
independently of the input;
\item \emph{robustness}: there is a computable positive rational radius
$\delta(c)$ such that
\[
U_c:=\{s\in S:d_S(s,J(c))<\delta(c)\}
\]
is nonempty, lies in $A$, has constant target membership, and, when
$\tau(c)$ is defined,
\[
F(U_c)\subseteq U_{\tau(c)}.
\]
\end{enumerate}
\end{definition}

Condition (E3) does not require an infinite amount of storage at any finite
step.  It requires a system with no fixed global counter cap across the
entire input family and all finite run prefixes.  This is precisely the
condition missing from a construction that supplies only a finite admissible
horizon.

\begin{theorem}[Conditional selection-reachability undecidability]
\label{thm:undecidable}
If a selected MTT realization satisfies
Definition~\ref{def:embedding}, then its selection reachability language is
undecidable.  More precisely,
\[
\operatorname{HALT}_{\mathcal M}
\leq_m
\operatorname{SEL}(\mathfrak C).
\]
\end{theorem}

\begin{proof}
Given a machine input $w$, compute $g(w)$ and hence the corresponding initial
MTT code $\iota(g(w))=J(c_0(w))$.  Conditions (E2) and (E3) identify every
finite MTT prefix with
the corresponding machine prefix.  By (E4), the MTT orbit reaches $H$ if and
only if $\mathcal M$ halts on $w$.  Therefore an algorithm deciding
$\operatorname{SEL}(\mathfrak C)$ would decide the halting language of the
fixed universal two-counter machine, a contradiction.  Conditions (E5) and
(E6) ensure that the reduction is not hidden in an input-dependent controller
and is stable on the declared neighborhoods.
\end{proof}

\begin{remark}[Current MTT status]
The theorem is exact and conditional.  Projection, non-Markovianity, stable
records, locality, or basin language alone does not construct $J$, $F$, the
zero test, an unbounded admissible register, or the neighborhoods $U_c$.
The current corpus has not supplied all six embedding clauses for a selected
physical MTT realization.
\end{remark}

\section{Finite Capacity Is a No-Go for the Strong Claim}

\begin{theorem}[Fixed-capacity reachability]\label{thm:finite}
Let $Q$ be finite and suppose each of two counters is bounded by a fixed
$C<\infty$.  For any deterministic transition rule on
\[
Q\times\{0,\ldots,C\}^2,
\]
halting reachability is decidable.
\end{theorem}

\begin{proof}
There are at most $|Q|(C+1)^2$ configurations.  Simulate from the initial
configuration.  Accept on halting.  If a configuration repeats before
halting, determinism makes the subsequent orbit periodic, so reject.
\end{proof}

\begin{corollary}[Capacity obstruction]
A selected realization with one fixed finite state carrier and computable
transition table cannot satisfy Definition~\ref{def:embedding}.  A family of
larger finite truncations is not enough unless the passage to arbitrarily
large admissible prefixes is itself uniform and proved.
\end{corollary}

\begin{proof}
Finite-state reachability is decidable by the same repeated-state argument,
whereas Theorem~\ref{thm:undecidable} would make the embedded halting language
undecidable.
\end{proof}

There are two coherent research targets:
\begin{enumerate}
\item prove a genuinely unbounded selected MTT carrier with robust operations
and no premature admissibility exit; or
\item retain a finite carrier and state only finite-horizon simulation and
verification results.
\end{enumerate}
The second target can still be physically and computationally valuable, but it
does not establish universal undecidability.

\section{Complexity and ``Full Simulation''}

\begin{definition}[Decision-time complexity]
Fix an encoding of inputs, a machine model, and a size function $|w|$.  A
decision-time lower bound for a language $L$ is a theorem that every deciding
algorithm in the declared model uses at least the stated resources on an
infinite family of inputs~\cite{HartmanisStearns}.
\end{definition}

\begin{remark}
Non-Markovianity says that a chosen reduced state is not sufficient for
autonomous one-step prediction.  Enlarging the state by memory variables may
restore Markovian evolution.  This fact neither proves undecidability nor a
complexity lower bound.
\end{remark}

\begin{remark}
Finite capacity usually makes a fixed model easier to decide by exhaustive
state exploration, though the state count can be large.  It does not by
itself imply computational irreducibility.
\end{remark}

The phrase ``prediction requires full simulation'' can mean at least three
different things:
\begin{enumerate}
\item no closed-form shortcut is presently known;
\item a particular prediction algorithm follows every step;
\item every correct algorithm requires resources comparable to stepwise
simulation.
\end{enumerate}
Only the third is a lower-bound claim, and it requires a formal computation
model and proof.  Theorem~\ref{thm:undecidable} establishes no such bound for
decidable instances and no claim about all MTT trajectories.

\section{Conditional Consequences for Control}

\begin{definition}[Universal intervention verifier]
Let $\mathcal U$ be an effectively presented family of admissible
interventions containing a no-intervention element $u_0$.  A universal
intervention verifier is a total algorithm $V(w,u)$ that returns ``yes'' if
the orbit obtained from input $w$ under intervention $u$ reaches $H$, and
``no'' otherwise.
\end{definition}

\begin{proposition}[Conditional intervention-verifier obstruction]
\label{prop:control}
On a realization satisfying Definition~\ref{def:embedding}, no universal
intervention verifier exists for a family containing the no-intervention
element $u_0$.
\end{proposition}

\begin{proof}
The map $w\mapsto V(g(w),u_0)$ would decide whether the uncontrolled encoded
orbit reaches $H$.  It would therefore decide
$\operatorname{HALT}_{\mathcal M}$, contradicting
Theorem~\ref{thm:undecidable}.
\end{proof}

This result does not forbid local feedback control, control on a decidable
subclass, probabilistic control, controller synthesis without a complete
verifier, or a verifier that sometimes returns ``unknown.''  Branching also
does not imply that interventions cannot alter branch weights or admissible
successors; that is a dynamical question.

\section{What A2 Does Not Derive}

\subsection{Probability}

An algorithmic no-go does not select a probability measure.  A probability
statement needs a measurable space, a normalized state or measure, and an
outcome map.  Undecidability is compatible with deterministic, stochastic,
or set-valued dynamics and therefore does not imply the Born rule.

\subsection{Quantization}

A finite survivor set may be discrete, but discreteness at one truncation is
not a quantization theorem.  Quantization requires a selected algebra,
representation, spectrum or integral condition, and a source relation to the
physical observables.

\subsection{Time, horizons, and irreversibility}

$N_{\max}$ is an iteration count.  It is not a physical time or a spacetime
horizon.  A noncompact ordering variable, a Lorentzian metric, and a
hyperbolic propagation law require their own bridge.  Likewise, a partial or
noninjective reduced map does not by itself prove thermodynamic
irreversibility.

\subsection{Upper dynamics}

Encoding termination says that the declared chart no longer supplies a
compatible representation.  It does not prove that upper evolution
terminates.  A different chart, extension, or upper description may remain
available.

\section{Scoped A2 Theorem}

\begin{theorem}[Conditional computability and prediction-depth package]
\label{thm:scope}
For a declared MTT chart-transition system:
\begin{enumerate}
\item prediction depth is a chart- and protocol-relative iteration diagnostic;
\item it is finite under the certified barrier hypothesis of
Theorem~\ref{thm:depth};
\item refinement monotonicity holds under the projection contract of
Proposition~\ref{prop:refinement};
\item finite-depth branching and nil boundaries represent underdetermination
or partial semantics, not algorithmic undecidability by themselves;
\item a robust uniform universal two-counter embedding implies undecidable
selection reachability by Theorem~\ref{thm:undecidable};
\item a fixed finite carrier has decidable reachability and cannot satisfy
that unbounded embedding contract; and
\item complexity, full-simulation, probability, quantization, physical time,
and universal control claims require additional hypotheses.
\end{enumerate}
\end{theorem}

\begin{proof}
Items 1--3 follow from the definitions and
Theorem~\ref{thm:depth}--Proposition~\ref{prop:refinement}.  Item~4 is
Propositions~\ref{prop:branch} and~\ref{prop:nil}.  Items~5 and~6 are
Theorems~\ref{thm:undecidable} and~\ref{thm:finite}.  Item~7 records the
missing structures identified in the preceding sections.
\end{proof}

\section{Version Delta and Research Frontier}

Relative to version~1, this revision:
\begin{itemize}
\item withdraws the unsupported no-global-predictive-map theorem;
\item replaces generic structural undecidability by a conditional many-one
reduction;
\item states the six-clause robust embedding contract explicitly;
\item corrects the conflict between unbounded counters and fixed finite
admissibility capacity;
\item distinguishes branching, undefined continuation, undecidability,
complexity lower bounds, and simulation cost;
\item narrows the control result to an intervention-verifier consequence of
the same reduction; and
\item removes the claimed derivation of quantization, measurement
probabilities, horizons, and irreversibility from prediction depth.
\end{itemize}

The decisive constructive frontier is now precise: exhibit, in one selected
MTT realization, a computable carrier and step map satisfying
(E1)--(E6), especially unbounded admissible persistence and a robust zero
test.  Until then, the corpus supports finite prediction-depth diagnostics
and finite simulations, but not a universal MTT undecidability theorem.

\section{Conclusion}

Finite admissibility can limit the range of a chosen description without
making its reachability language undecidable.  Branching can leave an outcome
unselected without making existential or universal reachability
uncomputable.  A nil boundary can make continuation undefined without
supplying a Boolean decision problem.  These distinctions do not weaken the
MTT program; they isolate the missing theorem.

If a selected realization robustly embeds a universal two-counter machine and
preserves every finite run prefix, undecidable selection reachability follows
immediately and rigorously.  If the realization instead has one fixed finite
capacity, reachability is decidable.  Program A2 therefore converts a broad
claim into a clean fork between an explicit universality construction and an
honest finite-horizon computational model.

\begin{thebibliography}{9}

\bibitem{A0}
P. Nero,
\emph{The Modal Triplet Theory Program A0: A Structural Theory of Reduced
Description}, version 2, 2026.

\bibitem{A1}
P. Nero,
\emph{The Modal Triplet Theory Program A1: Coherent Kinematics}, version 2,
2026.

\bibitem{Minsky}
M. Minsky,
\emph{Computation: Finite and Infinite Machines},
Prentice-Hall, 1967.

\bibitem{Turing}
A. M. Turing,
``On computable numbers, with an application to the
Entscheidungsproblem,''
\emph{Proceedings of the London Mathematical Society} 42 (1936), 230--265.

\bibitem{HartmanisStearns}
J. Hartmanis and R. E. Stearns,
``On the computational complexity of algorithms,''
\emph{Transactions of the American Mathematical Society} 117 (1965), 285--306.

\bibitem{Moore}
C. Moore,
``Unpredictability and undecidability in dynamical systems,''
\emph{Physical Review Letters} 64 (1990), 2354--2357.

\end{thebibliography}

\end{document}
